\begin{frame}{LIS}
    \metroset{block=fill}

    \begin{block}{Problema da maior subsequência crescente (LIS)}
        \footnotesize
        Dado um vetor com $n \leq 1000$ elementos, encontre sua maior subsequência crescente.
    \end{block}

    Exemplo: \pause
    \begin{center}
        \tikz[nodes={draw,circle}] {
            \node (1) at (0, 0) {3};
            \node (3) at (2, 0) {5};
            \node (6) at (5, 0) {2};
            \node (8) at (7, 0) {8};
            
            \only<2>{
                \node (2) at (1, 0) {2};
                \node (4) at (3, 0) {4};
                \node (5) at (4, 0) {6};
                \node (7) at (6, 0) {9}; 
            }
            \only<3>{
                \node[fill=red] (2) at (1, 0) {2};
                \node[fill=red] (4) at (3, 0) {4};
                \node[fill=red] (5) at (4, 0) {6};
                \node[fill=red] (7) at (6, 0) {9}; 
            }
        }
    \end{center} \pause

    Uma subsequência crescente com maior comprimento possível é $(2, 4, 6, 9)$ (em vermelho).

    Vamos definir uma função $\DP(i)$ como sendo a maior subsequência começando na posição $i$. Como definimos $\DP(i)$ recursivamente, isto é, a partir de $\DP(j)$ para $i < j$?

    \[ \DP(i) = \max_{i < j \text{ e } v_i < v_j}(\DP(j)+1) \ . \]

    % \begin{block}{Implementação da busca binária}
    %     \footnotesize
    %     \inputminted{cpp}{codigos/busca.binaria.cpp}    
    % \end{block}

\end{frame}
